% Generated draft: requires mathematical and source review.
\documentclass[11pt]{article}
\usepackage[margin=25mm]{geometry}
\usepackage{fontspec}
\setmainfont{Latin Modern Roman}
\usepackage{amsmath,amssymb,mathtools,unicode-math}
\setmathfont{Latin Modern Math}
\usepackage{enumitem,needspace,xurl,hyperref}
\setlength{\parindent}{0pt}
\setlength{\parskip}{5pt}
\setlength{\emergencystretch}{3em}
\newcommand{\N}{\mathbb{N}}
\newcommand{\Z}{\mathbb{Z}}
\newcommand{\Q}{\mathbb{Q}}
\newcommand{\R}{\mathbb{R}}
\newcommand{\C}{\mathbb{C}}
\newcommand{\F}{\mathbb{F}}
\newcommand{\A}{\mathbb{A}}
\newcommand{\PP}{\mathbb{P}}
\newcommand{\E}{\mathbb{E}}
\newcommand{\eps}{\varepsilon}
\newcommand{\dd}{\,\mathrm{d}}
\begin{document}
\section*{UnsolvedMath ID: 11000159}
\textbf{Source problem code:} catalogue entry
\subsection{Definitions and mathematical statement}
For a mapping class \(\varphi\) of a suitable finite-type surface, the Nielsen--Thurston classification distinguishes periodic, reducible nonperiodic, and pseudo-Anosov types. Periodic means finite order; reducibility involves preservation of an essential multicurve; the precise surface and exceptional-case conventions are not specified here. Retain \(d(\varphi)\) as the source's nonnegative distance, whose definition is absent.

\emph{Problem 2.2.} Determine how the Nielsen--Thurston type of \(\varphi\) is related to the four possibilities
\[
 d(\varphi)=0,\qquad d(\varphi)=1,\qquad
 d(\varphi)=2,\qquad d(\varphi)\geq3.
\]
The request is for the relationship between these classifications, including whatever implications or realizability restrictions hold in the intended setting. It does not assert a one-to-one correspondence between the three mapping-class types and the four distance cases.

The research summary suggests Hempel distance. In such an interpretation, one must first fix the Heegaard surface and disk sets \(\mathcal D_1,\mathcal D_2\) and the gluing convention, so that a possible distance is \(d_{\mathcal C(S)}(\mathcal D_1,\varphi(\mathcal D_2))\) in the curve graph. This is an explanatory conditional interpretation, not a definition supplied by the excerpt. It must not be replaced by the point displacement or stable translation length of \(\varphi\), which are different quantities. The source's trailing sentence says that the next three problems address the nonconstructive definition of \(d(\varphi)\); those problems and the definition are not included.
\subsection{Short English statement}
Relate the mapping class's periodic/reducible/pseudo-Anosov type to whether the source distance is \(0,1,2\), or at least \(3\). The distance and its auxiliary Heegaard data are missing, and a curve-graph translation length cannot be silently substituted.
\subsection{Sources}
\begin{itemize}
\item Local UnsolvedMath registry, permanent ID 11000159, source code catalogue entry, Problem 2.2, PDF page 156. No individual proposer is supplied; page distinguishes this label from repeated Problem 2.2 headings.
\item Farb, \emph{Problems on Mapping Class Groups and Related Topics}. Supplied URL: \url{https://www.math.uchicago.edu/~farb/papers/mcgbook.pdf}; recorded access 2026-07-29; extraction pdf-font-structure. The volume contains 329 explicitly typeset headings across contributed chapters.
\item The supplied research summary describes the excerpt as truncated and related to Hempel distance and mapping-class type, but gives no definition, theorem, or additional reference.
\item Supplied status open and OPEN-TRIAGE. The summary says it cannot verify the question. Current primary-literature status and redistribution terms for the publicly accessible author/publisher copy remain NEEDS\_REVIEW.
\end{itemize}
\end{document}
